\documentclass[10pt,a4paper]{amsart}
\usepackage[margin=19mm]{geometry}
\usepackage{amsmath,amssymb,mathtools}
\usepackage[scr=rsfs,cal=euler]{mathalfa}
\usepackage{xfrac}
\usepackage[unicode,hidelinks]{hyperref}
\newcommand{\ie}{i.e.\ }
\newcommand{\cf}{cf.\ }
\newcommand{\resp}{respectively\ }
\newcommand{\Subsection}{\subsection}
\providecommand{\nobreakdash}{\nobreak\hbox{-}}
\setlength{\emergencystretch}{2em}
\multlinegap0pt
\newcommand{\Changel}{\mathcode`l="7160}
\newcommand{\Changelback}{\mathcode`l="716C}
\newcommand{\NoChangel}[1]{%
\expandafter\let\csname old\string#1\endcsname=#1
\let#1=\relax
\newcommand{#1}{\mathop{\mathcode`l="716C\csname old\string#1\endcsname\mathcode`l="7160}\displaylimits}%
}

\NoChangel{\log}
\NoChangel{\ln}
\NoChangel{\lim}
\NoChangel{\limsup}
\NoChangel{\liminf}
\NoChangel{\varlimsup}
\NoChangel{\varliminf}
\NoChangel{\varinjlim}
\NoChangel{\varprojlim}

\newenvironment{smallpmatrix}{\big(\begin{smallmatrix}}{\end{smallmatrix}\big)}

\def\mto{\mathchoice{\longmapsto}{\mapsto}{\mapsto}{\mapsto}}

\let\<\langle
\let\>\rangle

\newcommand{\bG}{\boldsymbol{G}}
\newcommand{\SO}{\mathrm{SO}}

\newcommand{\Rrank}{\R\textup{-}\rank}

\newtheorem{theorem}{Theorem}
\newtheorem{lemma}{Lemma}
\newtheorem{corollary}[theorem]{Corollary}
\newtheorem{proposition}[theorem]{Proposition}
\newtheorem{problem}{Problem}
\newtheorem{question}[problem]{Question}
\theoremstyle{definition}
\newtheorem{conj}{Conjecture}
\newtheorem*{example}{Example}
\newtheorem{defn}{Definition}
\newtheorem{remark}{Remark}
\newtheorem{notation}{Notation}

\newcommand{\field}[1]{\mathbb{#1}}
\newcommand{\Q}{\field{Q}}
\newcommand{\R}{\field{R}}
\newcommand{\Z}{\field{Z}}
\newcommand{\C}{\field{C}}
\newcommand{\A}{\field{A}}
\newcommand{\F}{\field{F}}
\newcommand{\SL}{\mathrm{SL}}
\newcommand{\Sp}{\mathrm{Sp}}
\newcommand{\liesl}{{\frak{sl}}}
\newcommand{\rank}{{\operatorname{rank}}}

\let\D\Delta
\let\g\gamma
\let\G\Gamma
\let\l\lambda
\let\s\sigma

\newcommand{\frg}{\mathfrak g}
\newcommand{\fh}{\mathfrak h}

\let\bs\backslash

\title{Discrete linear groups containing arithmetic groups: Theorem 12}
\author{Indira Chatterji and T. N. Venkataramana}
\date{Source: J.\'E.P.--M. 12 (2025), 1605--1632; DOI 10.5802/jep.318}
\begin{document}\maketitle
\noindent\textit{Editorial scope.} The counted unit is original Theorem 12. The original notation and complete source-local measurable-map core used in its proof are retained as uncounted support. Later archimedean extensions, arithmetic applications and rank-one constructions are omitted. Original statement and proof wording, including the original notation in Notation 1, is preserved.
\section{Necessary original notation}
\begin{notation}\label{notation:Lie} Let $P$ be a minimal real parabolic subgroup of the simple Lie group~$G$ and denote by $N$ the unipotent in a maximal real split torus of $P$. Let $A$ be the connected component of identity in $S$. Denote by $P_0$ the subgroup $AN$ of $P$ (if~$G$ is split, then $P_0$ is the identity component of $P$; in general we have replaced the minimal parabolic subgroup $P$ by a subgroup $P_0$ which has no compact factors such that $P/P_0$ is compact). Let $K$ be a maximal compact subgroup of $G$. We~have the Iwasawa decomposition $G=KP_0=KAN$.
\end{notation}

\setcounter{theorem}{2}
\begin{theorem}[Main result]\label{theo:main}
Let $H$ be a semi-simple subgroup of a semi-simple Lie group $G$ with $\Rrank(H)\geq 2$ such that the normal subgroup of $G$ generated by~$H$ is all of $G$. Let $\Gamma $ be a Zariski dense subgroup of $G$ whose intersection with
$H$ is an irreducible lattice in $H$. Let $G=KAN$ be the Iwasawa
decomposition of $G$ and $P_0=AN$. If the isotropy of $H$ acting on $G/P_0$ is positive dimensional and non-compact at every point of $G/P_0$, then $\Gamma $ is a super-rigid subgroup of $G$.
\end{theorem}

An earlier version of Theorem \ref{theo:main} assumed $G$ to be simple. It was pointed out to us by Yehuda Shalom that the proof works for $G$ semi-simple as well.

\begin{theorem} \label{ranktwo} Let $\Gamma $ be a Zariski dense discrete subgroup of a simple Lie group~$G$ which intersects a semi-simple subgroup $H$ of $G$ (with $\Rrank(H)\geq 2$) in an irreducible lattice. Let $K$ be a maximal compact subgroup of $G$ and assume that $\dim (H)> \dim (K)$. Then $\Gamma$ is a super-rigid subgroup of $G$.
\end{theorem}

Here, \emph{super-rigid} is in the sense of Margulis \cite{M}. That is, all linear representations -- satisfying some mild conditions -- of the group $\Gamma$ virtually extend to (\ie coincide on a finite index
subgroup of $\Gamma$ with) a linear representation of the ambient group $G$.

\section{Preliminaries on measurable maps}

The aim of this section is to recall a few well known facts used in the proof of Theorem \ref{ranktwo} and prove Proposition \ref{rational} (this is a key fact in the proof of the main theorem).

\subsection{A mild generalisation of the Howe-Moore theorem}

An important ingredient in the proof is the following easy generalisation (Lemma \ref{mautner} below) of the ergodicity theorem of Moore (see \cite[Th.\,(2.2.6)]{Z}). This version is only slightly different from \cite[Th.\,(2.2.6)]{Z}; this is to take care of additional compact factors.

We give the proof for the sake of completeness. We~note that in the statement of Lemma \ref{mautner} below,
we may replace $\tilde{H}(\Z)$ by any irreducible lattice $\G$ in a real group~$H$ such that $\G$ maps densely
onto the maximal compact quotient of $H$; the arithmetic structure is not really used in the proof.

The almost $\Q$-simplicity of $\tilde H$ implies that there exists a connected absolutely almost simple simply connected group $H_0$ defined over a number field $K$, such that $\tilde H=R_{K/\Q}(H_0)$,
where $R_{K/\Q}(H_0)$ is the Weil restriction of $H_0$ from $K$ to $\Q$. Consequently,
\[
\tilde H (\R)=H_0(K\otimes \R)=\prod_{\alpha \in \infty} H_0(K_{\alpha}),
\]
where $\infty$ denote the set of equivalence
classes of archimedean embeddings of $K$. Let $A\subseteq\infty$ denote the archimedean embeddings $\alpha$ such that $H_0(K_{\alpha})$ is non-compact. Then, $\tilde H= H^*\times H^u$ where $H^u:=\prod
_{\alpha \in \infty \setminus A} H_0(K_{\alpha })$ is a compact group, and $H^*:= ~\prod_{\alpha \in A} H_0(K_{\alpha})$ is a semi-simple group without compact factors.

As in the case of the Moore ergodicity theorem of \cite{Z}, we~first show the vanishing of matrix coefficients.

\begin{lemma}\label{howemoore} Let $\tilde{H}$ be an almost $\Q$-simple, simply connected algebraic group with $\Rrank(\tilde{H})\geq 1$ and $\Delta < \tilde{H}(\Z)$ an arithmetic subgroup. Suppose that $\pi$ is a unitary representation of $\tilde H(\R)$ on a Hilbert space, such that for any \emph{non-compact} simple factor $H_{\alpha }$ of $H^*$, the space $\pi ^{H_{\alpha }}$ of vectors of $\pi $ invariant under the subgroup $H_{\alpha}$ is zero. Then, the space $\pi ^{S}$ of vectors in $\pi $ invariant under the non-compact subgroup $S<\tilde H(\R)$ is also zero.
\end{lemma}

\begin{proof} By the Howe-Moore theorem, (see \cite[Th.\,(2.2.20)]{Z}), for every pair of vectors $v,w\in \pi$ the matrix coefficient $\<g^*v,w\>$ tends to zero as $g^*$ tends to infinity in the group $H^*$.

Suppose $g(n)=(g^*(n),g_u(n)) \in \tilde{H}(R)=H^* \times H_u$ is a sequence which tends to infinity. Since $H_u$ is compact, we~may assume, by~passing to a subsequence, that $g_u(n)$ tends to an element $k\in H_u$.
Hence $g_u(n)v$ tends to some vector $v'\in \pi$. Then $g^*(n)$ also tends to infinity in $H^*$. Moreover, (write $g=g(n), g^*=g^*(n),g_u=g_u(n)$ for short):
\[
\<gv,w\>=\<g^*v',w\>+ \<g^*(g_uv-v'),w\>.
\]
The Cauchy Schwarz estimate and the unitarity of $g^*$ show that
\[
|\<g^*(g_uv-v'),w\>| \leq | g_uv-v'|\, |w|
\]
and hence tends to zero. By~the Howe-Moore theorem, $\<g^*v',w\>$ tends to zero. Hence the above implies that $\<gv,w\>$ tends to zero as $g= g(n)$ tends to infinity in $\tilde{H}(\R)$.

In particular, no non-compact subgroup of $\tilde H(\R)$ can fix a non-zero vector in $\pi$.
\end{proof}

\begin{lemma}\label{mautner} Let $\tilde{H}$ be an almost $\Q$-simple, simply connected algebraic group with
$\Rrank(\tilde{H})\geq 1$ and $\Delta <\tilde{H}(\Z)$ an irreducible lattice so that all the projections to the compact factors are dense. For any closed non-compact subgroup $S<\tilde{H}(\R)$, the group $\Delta$ acts ergodically on the quotient $\tilde{H}(\R)/S$.
\end{lemma}

\begin{proof} Let $V_0=L^2(\Delta \bs \tilde H(\R))$ be the space of square integrable functions on $\Delta \bs \tilde H(\R)$; the latter space has finite volume (\cite[Th.\,1]{BHC}),
and hence contains the space of constant functions. For any simple factor $H_{\alpha }$ of $H^*< \tilde H(\R)$, the space~$V_0^{H_{\alpha }}$ is just the space of constants, by~strong approximation (see \cite[Chap.\,II, Th.\,(6.7)]{M}; in the notation of \cite{M}, we~may take $B=\{{\alpha \}}$ to be a singleton). Consequently, if~$\pi $ denotes the space of functions in $V_0$ orthogonal to the constant functions, then $\pi $ satisfies the assumptions of Lemma \ref{howemoore}. Therefore, by~Lemma \ref{howemoore}, $\pi ^S=0$, and
hence the only functions on $\Delta \bs \tilde H(\R)$ invariant under the non-compact group $S$ are constants. This proves Lemma~\ref{mautner}.
\end{proof}

We now record a statement which will be used in the proof of Proposition \ref{rational}. We~thank the referee for pointing out the simple proof (and the correct formulation) of the following lemma.\enlargethispage{.5\baselineskip}

\begin{lemma}\label{useful} Let $\tilde{H}\mkern-3mu$ be a $\Q$-simple, simply connected algebraic group with \hbox{$\R$-rank$(\tilde{H})\mkern-5mu\geq\mkern-5mu 1$} and $\Delta <\tilde{H}(\Z)$ an arithmetic subgroup. Suppose that $s\in \tilde{H}(\R)$ generates an infinite discrete subgroup and let
$\tau:s^{\Z}\to \Z$ be an isomorphism. Then, there is no $s^{\Z}$\nobreakdash-equivariant Borel measurable map
\[
\phi^*:\Delta \bs \tilde{H}(\R)\to \Z.
\]
\end{lemma}

\begin{proof} The image (push-out) of the Haar measure on $\tilde{H}(\R)/\Delta$ under an $s^{\Z}$\nobreakdash-equi\-variant Borel measurable map $\phi^*$ gives a finite $\Z$-invariant measure on $\Z$, which cannot exist.
\end{proof}

\subsection{The Margulis super-rigidity theorem} The following is a version of the Margulis super-rigidity theorem, except that the Zariski closure of the image $\rho (\Delta)$ is not assumed to be an absolutely simple group and that $\rho (\Delta)$ is not assumed to have non-compact closure in $G'(k')$ if $k'$ is archimedean.


\setcounter{theorem}{8}
\begin{theorem}[Margulis]\label{super-rigid} Let $\tilde{H}$ be a $\Q$-simple simply connected algebraic group defined over $\Q$ of $\R$-rank~$(\tilde H)\geq 2$, let $\Delta < \tilde{H}(\Z)$ be a subgroup of finite index and $\rho : \Delta \to G'(k')$ a homomorphism into a linear algebraic group $G'$ over a local field $k'$ of characteristic zero.

\begin{enumerate}
\item If $k'$ is archimedean, then the map $\rho $ coincides, on a subgroup of finite
index, with a representation $\tilde \rho: \tilde H(\R) \to G'(k')$.

\item If the local
field $k'$ is non-archimedean, then $\rho (\Delta)$ is contained in a
compact subgroup of $G'(k')$.
\end{enumerate}
\end{theorem}

\begin{proof} The usual statement of Margulis' super-rigidity says that if $\rho$ is a homomorphism of an irreducible lattice $\Gamma $ in a real semi-simple linear Lie group $H$ without compact factors, and if $G'$ is an absolutely simple group defined over an archimedean local field $k'$, then a representation from $\Gamma $ into $G'(k')$ with Zariski dense image, extends to a smooth
representation of $H$ into $G'(k')$ (see \cite[Chap.\,VIII, Th.\,(C)]{M}). However, if $\tilde H$ is the group of
real points of a $\Q$-simple simply connected algebraic group, then $\tilde H(\R)$ may have compact factors and
hence $\tilde{H}(\Z)$ (or a finite index subgroup of $\tilde{H}(\Z)$) may have representations whose Zariski closures
have compact factors.

Even so, representations of $\tilde{H}(\Z)$ do extend to $\tilde H(\R)$ under the hypotheses of Theorem \ref{super-rigid}.
First of all, the Zariski closure $\mathcal H$ of any representation of a higher rank lattice is semi-simple \cite[Chap.\,VIII, Th.\,(3.10)]{M}. By~passing to a finite index subgroup of $\rho$ we may assume that this Zariski closure $\mathcal H$ is connected and semi-simple. We~may write $\mathcal H =H_1\cdots H_r$ as an almost direct product with each $H_i$ absolutely almost simple over an archimedean local field $k_i$ (which is $\R$ or $\C$). Hence the representation $\rho$ is of the form $\rho_1 \cdots \rho_r$, with each $\rho_i :\D \to H_i$ with Zariski dense image.

The $\Q$-simple group $\tilde{H}$ is obtained as the Weil restriction of scalars from $K$ to~$\Q$ of an absolutely almost simple simply connected group $H_0$ defined over a number field $K$. Namely, $\tilde{H}=R_{K/\Q}H_0$ (for the Weil restriction of scalar we refer to \cite[ Chap.\,6, p.\,115]{Z}). Then, $H(\R)=\prod_{\s \in K_\infty} H_0(K_\s)$ where $K_\infty$ is the set of inequivalent archimedean completions of the number field $K$.

Then, by~\cite[Chap.\,VIII, Th.\,(3.6)\,(ii)\,(a)]{M}, the representation $\rho_i$ coincides (up~to a homomorphism of $\D$ into the centre) with an algebraic representation of~$H_0$ into~$H_i$ defined over the field $k_i$ and a homomorphism $\s_i: K \to k_i$.
But any such archimedean embedding $\s_i$ is simply the inclusion of $K$ into $K_s$ (followed by a continuous embedding of fields $K_s \to k_i$) and hence the map $\rho_i :\D \to H_i(k_i)$ is the composite of the inclusion of $H_0(O_K)$ into $H_0(K_s)$ followed by an algebraic homomorphism $H_0 \to H_i$ defined over $k_i$.
In other words, $\rho_i$ is up to centre, an algebraic homomorphism of $H_0(K_s)$ into $H_i(k_i)$. Since this centre is finite, we~may pass to a further subgroup of finite index in $\D$ to ensure that this homomorphism $\rho $ restricted to the finite index subgroup, coincides with an algebraic representation of $\tilde H(\R)$.
\end{proof}

We can now prove the main result (Proposition \ref{rational}) of this section. In the proposition, when we talk of an algebraic subgroup $J$ of $G'(k')$ where $k'$ is an archimedean local field, we~mean that $J$ is an algebraic $\R$-subgroup of the \emph{real} group $R_{k'/\R}(G')$ obtained from $G'$ by the Weil restriction of scalars.

\begin{proposition}\label{rational} Let $\tilde{H}$ be an almost $\Q$-simple, simply connected algebraic
group with $\Rrank(\tilde{H})\geq 2$ and $\Delta <\tilde{H}(\Z)$ an arithmetic subgroup. Let $\rho :\Delta \to G'(k')$ be a representation, where $G'$ is a linear algebraic group over a local field $k'$ of characteristic zero. Suppose that $S<\tilde{H}(\R)$ a closed non-compact subgroup and let $J<G'(k')$ be an algebraic subgroup. Then any Borel measurable and $\Delta$-equivariant map $\phi :\tilde H(\R)/S\to G'(k')/J$ coincides with a rational
map almost everywhere on $\tilde{H}(\R)$. More precisely:
If $k'$ is an archimedean local field, there exist a homomorphism
$\tilde \rho : \tilde H(\R) \to G'(k')$ of real algebraic groups
defined over $\R$, and a point $p\in G'(k')/J$, such that for almost all $h\in \tilde{H}(\R)$, the map $\phi
(h)$ and the map $h\mto \tilde \rho (h)(p)$ coincide (that is to say, $\phi (h)=\tilde \rho (h)p$ for almost all $h\in \tilde H(\R)$ and coincides almost everywhere with an $\R$-rational map of real varieties on $\tilde{H}(\R)$).

If $k'$ is a non-archimedean local field, then the map $\phi$ is constant a.e.\ on $\tilde{H}(\R)$.
\end{proposition}

\begin{proof} Suppose first that $k'$ is archimedean. Let $\Delta '< \Delta $ be a subgroup of
finite index such that there exists (according to Theorem \ref{super-rigid} quoted above) a representation $\tilde \rho :\tilde
H(\R)\to G'(k')$ which coincides with $\rho $ on $\Delta '$. Consider the map $\phi ^*(h)=\tilde \rho (h)^{-1}(\phi (h))$ from $\tilde
H(\R)$ into the quotient $G'(k')/J$. Then, for all $\delta \in \Delta '$, almost all
$h\in \tilde H(\R)$ and all $s\in S$, we~have $\phi ^*(\delta h)= \phi ^*(h)$ and $\phi ^*(hs)=\tilde \rho (s)^{-1}(\phi ^*(h))$. That is,
the map $\phi ^*$ is $\Delta '$ invariant and $S$-equivariant for the action of $\tilde H(\R)$ on $G'(k')/J$ via the representation $\tilde
\rho $.\enlargethispage{-.5\baselineskip}%

The representation $\tilde \rho$ is algebraic; moreover, since $k'$ is
archimedean, by~assumption the group $J$ is a real algebraic subgroup
of $G'(k')$, and hence the action of $\tilde H(\R)$ on $G'(k')/J$ is
smooth. Let $S_1$ denote the Zariski closure of the image $\tilde
\rho (S)$. The $S_1$\nobreakdash-action on $G'(k')/J$ is smooth, hence the
quotient $S_1\bs G'(k')/J$ is countably separated. On~the other hand,
by Lemma \ref{mautner}, the action of $S$ on $\Delta '\bs \tilde
H(\R)$ is ergodic. Hence, by~\cite[Prop.\,(2.1.11)]{Z}, the
image of $\phi ^*$ is essentially contained in an $S_1$-orbit
\ie there exists a Borel set $E$ of measure zero in $\tilde H(\R)/S$,
such that the image under $\phi$ of the complement of $E$ is
contained in an $S_1$-orbit.

Since $S$ is a non-compact Lie group, $S$ contains an element $s$ of
infinite order which generates a discrete non-compact subgroup. We~may
replace $S$ by the Zariski closure of the group generated by the element $s$ and $S_1$
by the Zariski closure of the image of $s^{\Z}$. Hence the $S_1$-orbit of the preceding paragraph is of the
form $S_1/S_2$ with $S_2$ an algebraic subgroup of the abelian group
$S_1$.

\subsubsection*{Case 1}
Suppose that the inverse image $S'=S\cap \tilde \rho ^{-1}(S_2)$ is a
non-compact subgroup. By~Lemma \ref{mautner}, the group $S'$ acts
ergodically on $\Delta \bs \tilde H(\R)$; therefore, the map~$\phi ^*$ -- being $S'$ invariant -- is essentially constant: $\phi ^*(\Delta \bs \tilde
H(\R))=\{p\}$ for some point $p\in G'(k')/J$. That is $\phi (h)=\tilde
\rho (h)(p)$ a.e. on $\tilde{H}(\R)$, and coincides with a rational map a.e.

\subsubsection*{Case 2}
Suppose that $S'$ is compact. Since $s^{\Z}$ generates a discrete
non-compact subgroup and $S'$ is compact, the image $s_1^{\Z}:= \tilde
\rho (s^{\Z})$ also generates a discrete non-compact subgroup in
$S_1/S_2$ ($\tilde \rho $ being an algebraic, hence continuous, map).
Hence $s_1^{\Z}$-orbits in $S_1/S_2$ are closed so that the space
$s_1^{\Z}\bs S_1/S_2$ is countably separated. By~Lemma \ref{mautner},
$s^{\Z}$ acts ergodically on $\Delta \bs \tilde H(\R)$. By~applying
\cite[Prop.\,(2.1.11)]{Z} to the $s_1^{\Z}$-invariant map
$\overline{\phi ^*}: \Delta \bs \tilde H(\R)\to s_1^{\Z}\bs S_1/S_2$,
we deduce that the image of $\phi ^*$ is essentially contained in an
orbit of $s_1^{\Z}$ in the quotient group $S_1/S_2$. Then, by~Lemma~\ref{useful}, it~follows that $\phi ^*$ cannot exist and we are in Case 1.

If $k'$ is non-archimedean, then by Theorem \ref{super-rigid}, the
image $\rho (\Delta)$ is contained in a compact group $K$ which acts
smoothly on $G'(k')/J$ so that $K\bs G'(k')/J$ is countably
separated. The group $\Delta $ acts ergodically on $\tilde H(\R)/S$.
Hence, by~\cite[Prop.\,(2.1.11)]{Z}, the $S$-invariant map
$\overline{\phi} :\Delta \bs \tilde {H}(\R) \to K\bs G'(k')/J$ is
essentially constant, and therefore the image of $\phi$ is essentially
contained in an orbit of $K$.

Since $K$ is a compact subgroup of the $p$-adic group $G'(k')$, it~has
a decreasing sequence of open subgroups $(K_n)_{(n\geq 1)}$ (of
finite index in $K$) such that the intersection $\bigcap_{n\geq 1}
K_n=\{1\}$ is trivial. Then $\Delta_n=\Delta \cap \rho ^{-1}(K_n)$
is of finite index in $\Delta$. By~Theorem~\ref{super-rigid} applied to
$\Delta_n$, it~follows that the image of $\phi $ is contained in an
orbit of~$K_n$ for each $n\geq 1$. But since the subgroups $K_n$'s
converge to the identity subgroup, it~follows that the image of $\phi
$ is a singleton. That is, $\phi $ is constant on a co-null subset of $\tilde{H}(\R)/S$.
\end{proof}

\subsection{Some measure theoretic constructions}

We now mention a consequence of Fubini's theorem we will need in the proof of Theorem~\ref{ranktwo}.

\begin{notation} Suppose that $H$ is a locally compact Hausdorff
second countable topological group with a Haar measure $\mu$ and assume that
$(H,\mu)$ is $\sigma$-finite. Suppose that $(X,\nu)$ is a
$\sigma$-finite measure space on which $H$ acts such that the action
\hbox{$H\times X\!\to\! X$} --~denoted by $(h,x)\mto hx$ -- is measurable and so
that for each $h\in H$, the map $x\mto hx$ on $X$ preserves the
measure class of $\nu $. Let $Z$ be a measure space and let $f:X\to Z$ be
a measurable map.

If $\Sigma =\Sigma (X,Z)$ is the set of measurable maps from $X$ to $Z$, then $f\in \Sigma$ and $H$ acts on $\Sigma$ by left translations on $X$. We~regard two maps $\phi, \psi \in \Sigma$ to be equal if they coincide almost everywhere on $X$.

\end{notation}

\begin{lemma}\label{almostall} Under the above notation, suppose that there exists a co-null subset $X_1\subseteq X$ such that for each $x\in X_1$, the map $h\mto f(hx)$ is constant on a co-null subset $H_x$ of $H$.

Then, given $h\in H$, there exists a measurable subset $X_h$ which is co-null in $X_1$ and such that for all $x\in X_h$
\[
f(hx)= f(x).
\]

This equation says that $H$ lies in the isotropy of $f \in \Sigma (X,Z)$: for every $h\in H$, $hf=f$ in $\Sigma$.
\end{lemma}

\begin{proof} First, let $F=\{(x,a,b)\in X_1\times H \times H: f(ax)=f(bx) \}$. Being a pullback of the diagonal in $Z\times Z$ under a measurable map, the set $F$ is measurable. Further, for each $x$ in the co-null set $X_1$, the slice $\{x\}\times H_x \times H_x$ lies in $F$, and is co-null in $\{x\}\times H\times H$, so the set $F$ is co-null.

By Fubini (see \cite[Th.\,A, p.\,147]{Hal}) there exists $h_0\in H$ such that the intersection $E\times\{h_0\}:=F\cap (X\times H\times\{h_0\})$ is co-null in $X\times H\times\{h_0\}$, so that $E$ is co-null in $X\times H$.

Applying Fubini to $E$, there exists a co-null subset $H_1\subseteq H$ such that for each $h\in H_1$ there is $X_h\subseteq X$ co-null such that $f(hx)=f(h_0x)$ for all $x\in X_h$. Thus, the two $Z$-valued functions $x\mto f(hx)$ and $x\mto f(h_0x)$ are equal a.e. on $X$ and so lie in the same equivalence class of $\Sigma$.

On the set $\Sigma =\Sigma (X,Z)$ of measurable maps (modulo the equivalence that equality a.e. implies equivalent), the group $H$ operates by left translations on $X$. By~the preceding paragraph, $h f =h'f=h_0f$ for all $h,h' \in H_1$, so that $h^{-1}h'f=f$ lies in the isotropy of $f$.

Given $g\in H$, the set $H_1g$ is again co-null and so intersects $H_1$; that is, there exist $h,h'\in H_1$ such that $h'=hg$ and hence \emph{every} $g\in H$ is of the form $h^{-1}h'$ for some $h,h'\in H_1$. Hence, by~the preceding paragraph, $gf=f$ for all $g\in H$, proving the lemma.
\end{proof}

We recall a well known result of Furstenberg we will need, commonly known as Furstenberg's lemma but due to Zimmer in the form given below.

\begin{lemma}[{Furstenberg, \cite[Cor.\,(4.3.7) \& Prop.\,(4.3.9)]{Z}}]\label{furstenberg}
Suppose that $\Gamma $ is a closed subgroup of a locally compact topological group $G$ and that $P_0$ is
a closed amenable subgroup of $G$. Let $X$ be a compact metric $\Gamma $-space. Then there exists a Borel measurable $\Gamma $-equivariant map $\phi$ from a co-null subset of $G/P_0$ to $\mathcal{P}(X)$, the space of probability
measures on $X$.
\end{lemma}

\subsection{Some Lie theoretic results}

We recall that $G=G(\R)$ is a connected simple Lie group and $\G< G$ a Zariski dense discrete subgroup and $H<G$ a semi-simple subgroup (we denote by $H$ the group of real points $H(\R)$). In this section, we~suppose that $H$ is a semi-simple Lie subgroup of a \emph{simple} Lie group $G$ and refer to Notation~\ref{notation:Lie} of the introduction.
Let $\frg$ and $\fh$ be the Lie algebras of $G$ and $H$ respectively and $X=G/P_0$. Denote by $^{\g}(H)$ the conjugate $\g H \g ^{-1}$.

\begin{lemma} \label{G/Pgenerators} Let $G$ be a connected Lie group, $H$ a semisimple subgroup and $\Gamma<G$ a discrete Zariski dense subgroup. There exist finitely many elements $\g_1,\dots, \g_k \in \G$ such that $G$ is generated by the $H_i=\!{\ }^{\g_i}(H)$ for $i=1,\dots,k$. Moreover, for \emph{every} $x\in G/P_0$, the map $H_k\times \cdots \times H_1 \to G/P_0$ given by $(h_k,\dots,h_1)\mto h_k\cdots h_1x$, is~surjective.
\end{lemma}

\begin{proof} Let $U\subseteq G$ be a small symmetric (\ie $U=U^{-1}$) open neighbourhood of the identity. Then the countable union $\bigcup_{m=1}^{\infty} U^m$ is an open connected subgroup, so~that $G=\bigcup_{m=1}^{\infty} U^m$.

\Changel

Consider the sum $W=\sum_{\g \in \G} {\ }^{\g }(\fh) \subseteq \frg$ of the subspaces $^{\g}(\fh)$. This is stable under the action of $\G$ under the adjoint representation of $G$ and by the Zariski density of~$\G$, it~is stable under the adjoint action of $G$; the simplicity of $G$ implies that $\frg$ is irreducible for the action of $G$ and hence $W=\frg$. Hence there exist finitely many elements $\g_i \in \G$ such that
\[
\frg=\sum_{i=1}^l (^{\g_i}(\fh)).
\]
By the implicit function theorem, there exists a small symmetric open neighbourhood~$U$ of identity in $G$ such that $U$ is contained in the product set
\[
\Pi={\ }^{\g_l}(H)\cdots {\ }^{\g_1}(H).
\]
Hence the group generated by the product $\Pi$ contains $U^m$ for all $m$ and it follows that it is $G$.

Now, $G/P_0=\bigcup_{m=1}^{\infty} X_m$ with $X_m=U^m(P_0)$ an increasing sequence of open sets. The compactness of $G/P_0$ implies the existence of some $m$ such that $G/P_0=U^m(P_0)=X_m$. Thus, given $x\in G/P_0$, $x\in U^m(P_0)$; equivalently, $P_0=U^m(x)$. Hence $G/P_0=U^{2m}(x)$ for all $x\in G/P_0$.

By taking $U$ small enough, we~can assume that $U$ is contained in the $l$-fold product set $\Pi={}^{\g_1}(H)\cdots{} ^{\g_l}(H)$. It follows that $G/P_0 \subseteq (^{\g_1}(H)\cdots {}^{\g_l}(H))^{2m} (x)$ for \emph{every} $x\in G/P_0$:
\[
G/P_0=U^{2m}(x) \subseteq \Pi ^{2m}(x) \subseteq G/P_0.
\]
We may take $k=2ml$ and take the $(\g_j)_{1\leq j \leq k}$ to lie among the elements $\g_1,\dots, \g_l$ with possible repetitions. This proves the lemma.\Changelback
\end{proof}

\subsection{Invariance properties of a map associated to $\phi$}

In addition to the assumptions preceding Lemma \ref{G/Pgenerators} assume that the subgroup $H$ operates with non-compact isotropies on $G/P_0$. Suppose $\G$ is a Zariski dense discrete subgroup of $G$ whose intersection with $H$ is an irreducible lattice. Thus we are under the assumptions of Theorem~\ref{theo:main}.

\begin{lemma} \label{constantp} Let $H,G,\Gamma$ as in Theorem \ref{theo:main}. Suppose $Z$ is a countably separated topological space. Then, any $p: G/P_0 \to Z$ a measurable map which is $\G$-invariant, is~constant on a co-null subset of $G/P_0$.
\end{lemma}

\begin{proof} Consider the space $\Sigma $ of measurable maps from $G/P_0$ into $Z$. On~$\Sigma$ the group $G$ operates on the left by translations on $G/P_0$. Thus, for an element $g\in G$, $gp=p$ means that $p(g^{-1}x)=p(x)$ almost everywhere on $G/P_0$. The isotropy at the point $p$, of the action of $G$ on the set $\Sigma $, is~a subgroup of $G$. The invariance of the map $p$ under $\G$ implies that the isotropy at $p$ contains $\G$.

Suppose $X\subseteq G/P_0$ is the co-null subset on which $p$ is defined everywhere. Write~$H$ for $H(\R)$ as in Proposition \ref{rational}. For each $x\in X\subseteq G/P_0$ we have the map \hbox{$p_x: h\mto p(hx)$} from $H$ into $Z$. This map is defined almost everywhere on $H$. The invariance of~$p$ under $\Gamma $ implies that the map $p_x$ is invariant under $\D$. The isotropy~$H_x$ of~$H$ at~$x$ is non-compact. Since $Z$ is countably separated, the ergodicity of the action of $\D$ on $H/H_x$, then implies that the map $p_x$ is constant almost everywhere (by~\cite[Prop.\,(2.1.11)]{Z}). Then by Lemma \ref{almostall}, the isotropy of the map~$p$ contains~$H$. Hence it contains the groups $\g H \g ^{-1}$ for every $\g \in \G$. By~Lemma \ref{G/Pgenerators}, the isotropy is then all of $G$.

We then have, for each $g\in G$, $p(gx)=p(x)$ a.e. on $X$. Recall that as a topological space $G=K\times P_0$. In particular, for each $k\in K$, we~have $p(kx)=p(x)$ a.e. on $X$. Then by Lemma \ref{almostall}, there exists a co-null subset $Y\subseteq X$ such that for \emph{each} $x\in Y$, the equality $p(kx)=p(x)$ holds on a co-null subset $B \subseteq K$. The image of $B$ under the \emph{diffeomorphism} $k\mto kx$ from $K$ onto $G/P_0$ is therefore co-null and $p$ is constant on this image.
\end{proof}

Suppose we are under the hypotheses of Theorem \ref{theo:main}. Thus we have a Zariski dense discrete subgroup $\G$ of the group $G$, and $G$ contains the semi-simple subgroup $H$ whose intersection with $\G$ is an irreducible lattice $\D$
in $H$. We~have a homomorphism $\rho: \G \to G'(k')$ with Zariski dense image where $G'$ is a centreless absolutely simple group defined over a local field $k'$ and $P'<G'$ a minimal parabolic $k'$-subgroup. Consider the action of $G'(k')$ on the space $\mathcal P$ of probability measures on the compact Hausdorff space $G'(k')/P'(k')$. Then the quotient $Z=\mathcal{P}/G'(k')$ is a topological space (which is $(T_0)$ but is not necessarily Hausdorff). Then $\G$, which maps to $G'(k')$, acts on $\mathcal P$.

\begin{proposition} \label{oneorbit} Under the hypotheses of Theorem \ref{theo:main}, there exists a $\G$-equivariant measurable map $\phi: G/P_0 \to G'(k')/J$ for some closed subgroup $J< G'(k')$.
\end{proposition}

\begin{proof} By the Furstenberg lemma (Lemma \ref{furstenberg}), we~have a $\G$-\emph{equivariant} measurable map $\phi : G/P_0 \to \mathcal{P}$. Composing this with the quotient map $\mathcal{P} \to Z=\mathcal{P}/G'(k')$, we~have a $\G$-\emph{invariant} measurable map $p: G/P_0 \to Z$. By~\cite[Cor.\,(3.2.17)]{Z} (see also \cite[(2.1.9)]{Z}), the space $Z$ is countably separated. By~Lemma \ref{constantp} above, the map $p$ is a constant map, and hence the image of $\phi $ lies in an orbit of $G'(k')$ of a point~$\mu $ in~$\mathcal P$.

This orbit is isomorphic to the quotient $G'(k')/J$ for some closed subgroup $J$ of $G'(k')$ since orbits in $\mathcal P$ are locally closed \cite[Prop.\,(2.1.10) \& Cor.\,(3.2.17)]{Z}.
\end{proof}

\section{The non-archimedean result}
\Subsection{The non-archimedean case}

\begin{theorem}\label{nonarch} Suppose that $\Gamma $ is a Zariski dense discrete subgroup of a simple Lie group $G$ which intersects a semi-simple Lie subgroup $H$ (of real rank at least two) of~$G$ in an
irreducible lattice $\Delta$. Suppose that $H$ acts with non-compact isotropy at any point of $G/P_0$ {\rm(}or that $\dim (H)>\dim(K)$ for a maximal compact subgroup~$K$ of~$G$). Then, the group $\Gamma $ is
non-archimedean super-rigid (that is, if $G'$ is an absolutely almost simple group defined over a non-archimedean local field $k'$ of characteristic zero, then every representation $\rho : \G \to G'(k')$ has compact image{\rm)}.
\end{theorem}

\begin{proof} Suppose that $\rho :\Gamma \to G'$ is a representation of $\Gamma $ into an absolutely simple (centreless) algebraic group $G'$ defined over a \emph{non-archimedean} local field $k'$ of characteristic zero with Zariski dense image.

By Proposition \ref{oneorbit}, there exists a closed subgroup $J< G'(k')$ and $\G$-equivariant measurable map $\phi : G/P_0\to G'(k')/J$. By~Proposition \ref{rational}, for any $x\in G/P_0$, the map $\phi_x: h\mto \phi(hx)$ is constant a.e. in $H$ (we are therefore using the non-archimedean version of Margulis' super-rigidity as in Theorem \ref{super-rigid} to conclude that the
image of the lattice in $H$ is contained in a compact subgroup in $G'(k')$).

Let $\Sigma $ now denote the set of measurable maps from $G/P_0$ into $Z=G'(k')/J$. Then~$\phi $ lies in $\Sigma$. The constancy of the map $\phi_x$ for almost all $x\in G/P_0$ then implies as in the proof of Lemma \ref{constantp}, that $h\phi =\phi$ for all $h\in H$. Thus $H$ leaves $\phi $ invariant. The equivariance of $\phi $ under $\G$ and invariance under $H$ then implies that for every $\g \in \G$, the conjugate group $^{\g} (H)$ also leaves $\phi $ invariant and hence the isotropy of~$\phi$ contains the group generated by these conjugates. By~Lemma \ref{G/Pgenerators}, this group is precisely~$G$. Hence $G$ leaves $\phi $ fixed. By~Lemma \ref{constantp}, the map $\phi $ is then constant. Hence $\rho (\Gamma)$ fixes a point $\mu $ in ${\mathcal P}$.

But the isotropy subgroup $J$ in $G'(k')$ of a probability measure $\mu\in {\mathcal P}$ is (by~\cite[Cor.\,(3.2.19)]{Z})
either compact, whence $\rho (\G)$ is contained in a compact group, or else the isotropy is contained in an algebraic group $L < G'$ of strictly smaller dimension. The latter is impossible because we have assumed that $\rho (\G)$ is Zariski dense in $G'$. Therefore, $\rho (\G)$ lies in a compact subgroup of $G'(k')$.

This means that $\Gamma $ is non-archimedean super-rigid in $G$, and proves Theorem \ref{nonarch}.
\end{proof}

\subsection{Original dimension criterion used in the alternative hypothesis}
\setcounter{lemma}{10}
\begin{lemma}\label{trick} Let $H$ be a semi-simple subgroup of simple
group $G$ and $K$ a maximal compact subgroup of $G$. Assume that
$\dim(H)>\dim(K)$. Let $G=KAN$ be an Iwasawa decomposition of $G$ and
$P_0=AN$. Then the isotropy subgroup of $H$ at any point in $G/P_0$ is
a non-compact subgroup of $H$.
\end{lemma}

\begin{proof} Since $G/P_0=K$, we~have $\dim(G/P_0)=\dim(K)$, and since
$\dim(H)>\dim(K)$, at any point $p\in G/P_0$, the isotropy of $H$ is a
positive dimensional subgroup, which is conjugate to a subgroup of
$P_0$; the latter has no compact subgroups, hence the isotropy of $H$
at $p$ is a non-compact subgroup.
\end{proof}


\begin{thebibliography}{Zim84}
\bibitem[BHC61]{BHC} A. Borel and Harish-Chandra, ``Arithmetic subgroups of algebraic groups'', \emph{Bull. Amer. Math. Soc.} \textbf{67} (1961), 579--583.
\bibitem[Hal74]{Hal} P. R. Halmos, \emph{Measure theory}, Graduate Texts in Math. 18, Springer, Cham, 1974.
\bibitem[Mar91]{M} G. A. Margulis, \emph{Discrete subgroups of semisimple Lie groups}, Ergeb. Math. Grenzgeb. (3) 17, Springer-Verlag, Berlin, 1991.
\bibitem[Zim84]{Z} R. J. Zimmer, \emph{Ergodic theory and semisimple groups}, Monographs in Math. 81, Birkh\"auser Verlag, Basel, 1984.
\end{thebibliography}
\end{document}
